\chapter*{Introduction Générale}
\addcontentsline{toc}{chapter}{Introduction Générale}
\markboth{Introduction Générale}{Introduction Générale}

L'essor considérable des technologies numériques au cours des deux dernières décennies a engendré une production massive de données dans pratiquement tous les secteurs d'activité~: santé, agriculture, finance, industrie, transport et administration publique. Cette accumulation exponentielle d'informations, souvent désignée sous le terme de \textit{Big Data}, constitue à la fois une opportunité et un défi. Une opportunité, car ces données recèlent des connaissances précieuses susceptibles d'éclairer la prise de décision~; un défi, car leur volume, leur variété et leur vélocité dépassent les capacités des méthodes d'analyse traditionnelles.

C'est dans ce contexte que la \textbf{science des données} s'est imposée comme une discipline à part entière, à l'intersection de l'informatique, des statistiques et de l'expertise métier. Elle propose un ensemble de méthodes et d'outils permettant d'extraire de la valeur à partir de données brutes, en passant par des étapes structurées de collecte, de nettoyage, d'exploration, de modélisation et de restitution des résultats. Parmi les outils les plus puissants dont dispose cette discipline figurent les \textbf{modèles d'intelligence artificielle (IA)}, qui offrent la capacité de détecter automatiquement des régularités, d'effectuer des prédictions et de soutenir des processus décisionnels complexes.

Ce rapport s'inscrit dans le cadre du module de Science des Données du \textbf{Master 2 en Informatique, option Intelligence Artificielle}, à l'Université de Dschang. Il a pour objectif de présenter, de manière structurée et rigoureuse, les fondements théoriques et méthodologiques de l'exploitation des données à l'aide des modèles d'IA, puis de les illustrer à travers une application concrète~: la \textbf{prédiction du rendement agricole} pour faciliter la prise de décision, concrétisée sous forme de la plateforme \textbf{AgriYield}.

Le choix de cette application n'est pas anodin. L'agriculture constitue un pilier économique fondamental pour de nombreux pays, en particulier en Afrique subsaharienne, où elle emploie une part importante de la population active et contribue significativement au produit intérieur brut. La capacité à anticiper les rendements agricoles à partir de données climatiques, agronomiques et géographiques représente un levier stratégique pour la sécurité alimentaire, l'allocation optimale des ressources et la planification des politiques publiques.

\bigskip
\noindent Le présent rapport est organisé autour de \textbf{cinq chapitres complémentaires}~:

\begin{itemize}[leftmargin=1.5cm]
  \item Le \textbf{Chapitre I} pose les fondements théoriques de l'exploitation des données, en clarifiant les notions essentielles relatives aux types de données, à leur cycle de vie et aux exigences de qualité qui conditionnent la fiabilité des analyses.
  \item Le \textbf{Chapitre II} présente les grandes familles de modèles d'intelligence artificielle mobilisables dans un projet de science des données, en distinguant les approches supervisées, non supervisées et les réseaux de neurones profonds.
  \item Le \textbf{Chapitre III} décrit la méthodologie opérationnelle d'exploitation des données avec les modèles d'IA, depuis la formulation du problème jusqu'à l'évaluation et l'interprétation des résultats.
  \item Le \textbf{Chapitre IV} constitue l'application pratique du rapport~: il mobilise les concepts et méthodes des chapitres précédents pour construire un système de prédiction du rendement agricole à partir d'un jeu de données réel (FAO/Kaggle).
  \item Le \textbf{Chapitre V} propose une prise de recul critique sur les apports, les limites et les perspectives des modèles d'IA dans l'exploitation des données.
\end{itemize}

\bigskip
À travers cette progression, le rapport vise à offrir une vision à la fois théorique et appliquée de la science des données, en montrant comment les modèles d'IA peuvent être concrètement mis au service de problématiques réelles à fort impact sociétal.

\newpage
